\documentclass{article}
\begin{document}
\section{Introduction}
Bubble nucleation starts boiling. It happens at the nanoscale, so experiments struggle to resolve it. Molecular dynamics fills this gap. In this work, we use it to study how surface patterns affect nucleation on a heated wall.

\section{Simulation method}
We model argon on a platinum wall. The wall carries a periodic pattern of hydrophilic and hydrophobic stripes. We vary the stripe width from 2 nm to 8 nm. We heat the wall suddenly and track the vapour phase with a local density criterion.

\section{Results}
Figure 2 shows snapshots of the liquid near the wall. Vapour first appears above the hydrophobic stripes. It then spreads sideways until it meets the hydrophilic stripes. We find that the hydrophilic stripes pin the bubble edge.

We now look at the onset time. Figure 3 shows it against stripe width. Narrow stripes give a short onset time. Wide stripes behave almost like a uniform hydrophobic wall. The trend is clear, but the scatter between runs is large.

The energy barrier explains part of this trend. On hydrophobic patches the barrier is lower, so vapour forms first there. The hydrophilic stripes then hold the bubble in place. We think this pinning keeps the bubble small and stable.

We also checked the effect of wall temperature. Higher temperatures shorten the onset time on all patterns. However, the ranking of patterns stays the same. Pattern width therefore matters across the range we studied.

\section{Conclusions}
In this work, we showed that striped walls control where bubbles form. Vapour appears on hydrophobic stripes and is pinned by hydrophilic ones. These results suggest a route to stable, well-placed nucleation sites.
\end{document}
